\section{Transfer through a rectangle and the number of crossings}
\label{sec:transfer}

Our first geometric estimate bounds the number of times the double-dimer
loops can cross a slab of fixed positive width.  We prove it in a finite
Temperleyan rectangle.  A particle representation makes a large number of
crossings expensive: changing the matching assignments at selected cap
arcs forces the transfer into particle sectors with a quadratic spectral
penalty.  We include the boundary calculation because the deleted corner
is part of the measure under consideration.

\subsection{Rectangles and traversals}

After translation, rotation, and possibly reflection, the rectangle has
vertices
\[
 R_\delta=\{(\delta x,\delta y):1\le x\le m,\ 1\le y\le n\}
                 \setminus\{(\delta,\delta)\},
 \qquad n=2p+1,
\]
where both $m$ and $n$ are odd.  Nearest-neighbor edges join the remaining
vertices.  We consider sequences for which $\delta m$ and $\delta n$
converge to positive finite limits.  The deleted corner is on the first
column.  We transfer in the horizontal direction and call it the
longitudinal direction.  Geometric cut locations use physical coordinates;
transfer calculations divide them by $\delta$ and count columns.
A seam is a line between two consecutive columns.
The part of the rectangle on either side of a seam is called a
\emph{cap}.

Two path portions are called \emph{parameter-disjoint} when their parameter
interiors are disjoint; they may share endpoints. For two longitudinal
cuts $u<v$, a loop that does not visit both $\{x\le u\}$ and
$\{x\ge v\}$ contributes zero traversals. For every other loop,
follow its orientation and record its successive visits to
$\{x\le u\}$ and $\{x\ge v\}$, suppressing repeated visits to the same
side before the other side is visited.  The resulting cyclic word alternates and has even length.  Its length is the number
of \emph{traversals} of the slab by that loop.  Sum over all loops to
obtain $N(u,v)$.
Equivalently this is the maximal number of parameter-disjoint trips
between the two cuts, in either direction.  In particular any chosen
family of disjoint crossing subarcs is bounded by $N(u,v)$, even if
those subarcs have excursions elsewhere.  Every traversal crosses any
intermediate seam, so
\begin{equation}\label{eq:seam-max-crossings}
 N(u,v)\le n.
\end{equation}
All cuts below may be moved by one mesh step to seams.

\begin{proposition}[Crossings in a finite rectangle]
\label{prop:rectangle-crossings}
Let the odd corner-rooted rectangles $R_\delta$ converge to a fixed
nondegenerate rectangle.  Let $u_\delta<v_\delta$ be parallel cuts in
either lattice direction, whose limiting coordinates are strictly
inside the longitudinal sides and whose limiting separation is
positive.  Under two independent uniform matchings of $R_\delta$,
\[
 \lim_{K\to\infty}\limsup_{\delta\downarrow0}
      \PP\{N(u_\delta,v_\delta)>K\}=0.
\]
The traversals are counted over all loops, with no restriction on their
portions outside the slab.
\end{proposition}

The proof will compare the ordinary partition function with modified
counts obtained by changing the allowed matching assignments in the caps.
We first develop the column transfer and its boundary normalization.
The same calculation also gives the uniform summability of spectral
states needed for the short-interval estimate in
Section~\ref{sec:short-segment}.

\subsection{Particles and column transfer}

Take black sites in the first column to have odd row indices.
A seam's intersection with a horizontal dimer is described by an
occupancy bit $b_i\in\{0,1\}$ in row $i$.
On a seam, define its particle bit $q_i$ to be $b_i$ if the site immediately
to its left is black, and $1-b_i$ otherwise.  At the initial boundary we
use the checkerboard signs of a virtual column $0$ and a single occupied
incoming edge at the deleted corner.  This edge simply declares that
corner already matched; it is not an edge of the rectangle.  The initial
particle set is therefore
\[
 \alpha=\{3,5,\ldots,2p+1\},
\]
and, because $m$ is odd, the empty final boundary has particle set
\[
 \beta=\{2,4,\ldots,2p\}.
\]
Both sets have size $p$.

Consider a column whose black rows are odd.  At a black row the incoming
and outgoing particle bits are $1-b_i^{\rm in}$ and $b_i^{\rm out}$;
at a white row they are $b_i^{\rm in}$ and $1-b_i^{\rm out}$.
An incoming or outgoing horizontal dimer leaves the particle unchanged.
A vertical dimer instead moves one particle from its black endpoint to
its white endpoint.  Thus particles stay or hop to an adjacent row, with
hops allowed only from black to white rows.  Disjoint such moves determine
the dimers in the column uniquely: the rows unused by horizontal dimers
are paired by the hops.

Let $A$ be the $n\times n$ matrix whose rows index outputs and columns
index inputs:
\begin{equation}\label{eq:one-particle-transfer}
 A_{ij}=\1_{\{i=j\}}
       +\1_{\{j\ \mathrm{odd},\ |i-j|=1\}}.
\end{equation}
In the sector with $j$ particles, the transfer is $\Lambda^j A$, written
in the increasing-index wedge basis. This is the nonintersecting-path
determinant viewpoint of Lindstr\"om and Gessel--Viennot
\cite{Lindstrom,GV}; in this one-column setting its positivity can be
checked directly. To check its signs, an allowed
bijection between source and destination rows cannot reverse their order:
such a reversal would require opposite hops across the same adjacent
pair, whereas only black-to-white hops are allowed.  A hop also cannot
pass a staying particle.  Hence every surviving determinant term has
positive sign and is exactly one legal collection of disjoint moves.
The transfers in consecutive columns alternate between $A$ and $A^*$.
For two independent matchings we take their tensor product.

Write $T^{(j)}_{b,a}$ for the transfer in sector $j$ from the seam after
column $a$ to the seam after column $b$, so it contains $b-a$ column
steps.  The number of single matchings is
\[
 Z=\langle e_\beta,T^{(p)}_{m,0}e_\alpha\rangle,
 \qquad Z_{\rm dd}=Z^2
\]
for the two-replica partition function.  All these matrix entries count
matchings with positive weights.

\subsection{Singular modes and the cost of changing particle number}

The alternating transfers have compatible singular bases. The
open-boundary sine modes belong to the transfer-matrix analysis developed
by Lieb~\cite{Lieb} and Rasmussen--Ruelle~\cite[Section~IV~A]{RR}.
We give the derivation with the alternating bases and odd-width zero mode
needed by the cap calculation. This gives
both the large-sector penalty used below and the summability needed for
the short-interval estimate in the next section.

\begin{lemma}[Spectrum and propagation]\label{lem:spectrum}
For $n=2p+1$, put
\[
 u_l=\operatorname{arsinh}\!\left(\cos\frac{\pi l}{n+1}\right),
 \quad 1\le l\le p,
 \qquad s=\exp\!\left(\sum_{l=1}^p u_l\right).
\]
There are orthonormal one-particle bases on alternate seams in which
each transfer is diagonal with entries
\[
 e^{u_1},\ldots,e^{u_p},1,e^{-u_p},\ldots,e^{-u_1}.
\]
Every basis vector has site coordinates bounded by $C/\sqrt n$, for an
absolute constant $C$.  Moreover,
\begin{equation}\label{eq:mode-gap}
 u_l\ge c\frac{p+1-l}{n}
\end{equation}
for an absolute $c>0$.  If $t=b-a$ and $1\le k\le p$, then
\begin{equation}\label{eq:sector-penalty}
 \frac{\|T^{(p+k)}_{b,a}\|\,
             \|T^{(p-k)}_{b,a}\|}{s^{2t}}
 \le \exp\!\left(-c\frac{t k^2}{n}\right).
\end{equation}
Finally, assign to any subset $J$ of the one-particle modes the normalized
propagation weight
\[
 w_t(J)=s^{-t}\prod_{j\in J}\sigma_j^t,
\]
where the $\sigma_j$ are the displayed singular values.  For every
$\eta>0$,
\begin{equation}\label{eq:all-state-sum}
 \sum_{J\subseteq\{1,\ldots,n\}}w_t(J)
 =2\prod_{l=1}^p(1+e^{-t u_l})^2\le C_\eta,
 \qquad t\ge\eta n.
\end{equation}
The corresponding sum for two replicas is at most $C_\eta^2$.
\end{lemma}

\begin{proof}
Split the site space into its odd and even rows.  The hop block from odd
to even rows has singular values
\[
 \lambda_l=2\cos\frac{\pi l}{2p+2},\qquad 1\le l\le p,
\]
and a one-dimensional kernel on the odd rows.  For completeness, on the
even rows its square is the tridiagonal matrix with diagonal $2$ and
adjacent entries $1$; its normalized eigenvectors have coordinates
proportional to $\sin(\pi l r/(p+1))$, $1\le r\le p$.
The associated normalized odd modes are proportional to
$\sin(\pi l(2r-1)/(2p+2))$, $1\le r\le p+1$.
Both families have coordinates bounded by $C/\sqrt n$.
The odd kernel vector has alternating coordinates of magnitude
$1/\sqrt{p+1}$.

On the odd/even pair belonging to $\lambda_l$, the matrix $A$ is
\[
 \begin{pmatrix}1&0\\ \lambda_l&1\end{pmatrix}.
\]
Its singular values are $e^{u_l},e^{-u_l}$ because
$\lambda_l=2\sinh u_l$.  Its expanding right and left vectors are
$(a_l,b_l)$ and $(b_l,a_l)$, respectively, where
\begin{equation}\label{eq:mode-components}
 a_l^2=\frac{e^{u_l}}{2\cosh u_l},\qquad
 b_l^2=\frac{e^{-u_l}}{2\cosh u_l},\qquad
 a_l\ge 2^{-1/2}.
\end{equation}
The contracting vectors may be taken as $(-b_l,a_l)$ and $(-a_l,b_l)$.
These two-dimensional changes of basis preserve the site bound.
Writing $A=U\Sigma V^*$, an $A$ step sends the $V$ basis to the $U$ basis,
and the next $A^*$ step sends the $U$ basis back to the $V$ basis, with
the same positive diagonal $\Sigma$.

For $r=p+1-l$ we have
$\cos(\pi l/(2p+2))=\sin(\pi r/(2p+2))$.
The elementary lower bounds for sine on $[0,\pi/2]$ and for
$\operatorname{arsinh}$ on $[0,1]$ give \eqref{eq:mode-gap}.
The largest product in sector $p$ fills all expanding modes and equals
$s$.  In sector $p+k$ one additionally fills the zero mode and the
$k-1$ least contracting modes; in sector $p-k$ one removes the $k$
weakest expanding modes.  Consequently the left side of
\eqref{eq:sector-penalty} is exactly
\[
 \exp\!\left[-t\left(
       \sum_{r=1}^{k}u_{p+1-r}
       +\sum_{r=1}^{k-1}u_{p+1-r}\right)\right],
\]
which proves the bound, including $k=1$.

Relative to the filled expanding modes, an arbitrary subset is specified
by holes in expanding modes, occupied contracting modes, and the free
choice of the zero mode.  Summing their weights gives the product in
\eqref{eq:all-state-sum}.  The logarithm of the full sum is bounded by
$\log 2+2\sum_{r\ge1}e^{-c\eta r}$, by \eqref{eq:mode-gap}.
\end{proof}

\subsection{The boundary caps}

A spectral estimate for the middle of the rectangle is useful only if
the boundary factors have the correct normalization.  Let
\[
 L_a=T^{(p)}_{a,0}e_\alpha,
 \qquad R_b=(T^{(p)}_{m,b})^*e_\beta
\]
be the ordinary single-replica cap vectors.  The double-replica vectors
are $L_a\otimes L_a$ and $R_b\otimes R_b$.

\begin{lemma}[Cap normalization]\label{lem:cap-comparison}
Fix $\eta,\rho>0$ and integer seams $0\le a\le b\le m$.
If $m/n\ge\rho$, $a\ge\eta n$, and $m-b\ge\eta n$, then
\begin{equation}\label{eq:cap-comparison}
 \frac{\|L_a\otimes L_a\|\,
             \|R_b\otimes R_b\|\,s^{2(b-a)}}{Z_{\rm dd}}
 \le C_{\eta,\rho}.
\end{equation}
In particular the bound is uniform as the two cuts vary in any fixed
longitudinal interior region of a nondegenerate limiting rectangle.
\end{lemma}

\begin{proof}
The initial wedge occupies all odd sites except the corner.  In the
orthonormal odd-mode basis it is a sum of wedges with one hole.  If $h_0$
is the coefficient for a hole in the zero mode, and $h_l$ that for a
hole in paired mode $l$, then the corner coordinates of these modes give
\begin{equation}\label{eq:corner-hole}
 |h_0|=\frac1{\sqrt{p+1}},\qquad
 \frac{|h_l|^2}{|h_0|^2}
 =2\sin^2\frac{\pi l}{n+1}\le2.
\end{equation}
Every term except the zero-mode hole has an occupied zero mode, which
can never reach the all-even final wedge.  Thus only the zero-mode-hole
term contributes to $Z$.

Since $m$ is odd, the odd-to-even transfer in pair $l$ is
\[
 a_l^2e^{m u_l}-b_l^2e^{-m u_l}
 =a_l^2e^{m u_l}\bigl(1-e^{-2(m+1)u_l}\bigr).
\]
With the overall wedge signs fixed by the positive partition function,
we obtain the exact formula
\begin{equation}\label{eq:rectangle-partition}
 Z=|h_0|\left(\prod_l a_l^2\right)s^m
       \prod_l\bigl(1-e^{-2(m+1)u_l}\bigr).
\end{equation}
There is no cancellation between different hole terms, as all the other
terms have zero output coefficient.  By \eqref{eq:mode-gap} and $m/n\ge\rho$,
\begin{equation}\label{eq:partition-product}
 0<c_\rho\le
 \prod_l\bigl(1-e^{-2(m+1)u_l}\bigr)\le1.
\end{equation}
Indeed this product is bounded below by a fixed positive product
$\prod_{r\ge1}(1-e^{-c\rho r})$.

We next estimate the input cap without discarding its other hole terms.
Put
\[
 f_l(q)=a_l^2e^{2q u_l}+b_l^2e^{-2q u_l}.
\]
Different hole terms remain orthogonal after propagation: in each pair
the number of particles is preserved, as is the zero-mode occupancy.
It follows that
\begin{equation}\label{eq:left-cap-exact}
 \|L_a\|^2
 =|h_0|^2\prod_l f_l(a)
       \left(1+\sum_l\frac{|h_l/h_0|^2}{f_l(a)}\right).
\end{equation}
The product divided by
$\prod_l a_l^2 e^{2a u_l}$ is
$\prod_l(1+(b_l/a_l)^2e^{-4a u_l})$, and the sum in parentheses is at
most $1+C\sum_l e^{-2a u_l}$.  Both are uniformly bounded when
$a\ge\eta n$.  Therefore
\[
 \|L_a\|\le C_\eta |h_0|\left(\prod_l a_l\right)s^a.
\]
The all-even final wedge has one particle in every paired mode and no
zero-mode particle.  The same calculation, now without hole terms,
gives
\[
 \|R_b\|^2=\prod_l f_l(m-b),\qquad
 \|R_b\|\le C_\eta\left(\prod_l a_l\right)s^{m-b}.
\]
Multiplying these bounds, inserting $s^{b-a}$, and comparing with
\eqref{eq:rectangle-partition}--\eqref{eq:partition-product} proves the
single-replica comparison.  Its square is \eqref{eq:cap-comparison}.
\end{proof}

\subsection{Cap diagrams and modified endpoint signs}

We now turn the spectral penalty into a geometric estimate.  Superpose
the two matchings inside one cap, retaining the multiplicities of its
edges but forgetting which matching supplies each single edge.  After
doubled edges are suppressed, each component is either an interior
closed loop or a simple arch joining two single seam edges.  The arches
are disjoint.  An ordinary interior loop has two alternating matching
assignments; a doubled edge has only one.

At a single seam edge, the difference of the two particle bits is $+1$
or $-1$.  It equals the signed crossing in the direction of the cycle
obtained by directing first-matching edges from black to white and
second-matching edges from white to black.  Each arch has opposite
particle-difference signs at its two ends.  Either choice of those
opposite signs determines its alternating matching assignment uniquely.
These descriptions also hold for the cap at the missing corner: the
virtual incoming edge occurs in both replicas and produces no open
single strand.

A \emph{tear} is the following formal modification of an arch.  Instead
of requiring opposite endpoint signs, prescribe both signs to be $+$,
or both to be $-$, with weight one for that prescription.  The uncolored
arch and all its edge multiplicities remain unchanged.  This defines a
vector of seam states even though the modified assignment need not be
a pair of matchings inside the cap.  A positive tear increases the
first particle number by one and decreases the second by one; a negative
tear has the reverse effect.  Geometrically it can be regarded as a
source or a sink on the arch.

Here is the coefficient rule for the resulting formal vector.  Fix an
uncolored cap diagram $D$ and its prescribed tears, and let $c(D)$ be
its number of internal single loops.  At a seam site with multiplicity
zero or two, the two occupancy bits are forced, and hence so are the
particle bits.  At the single seam sites, choose endpoint signs subject
to the opposite-sign rule on every unmarked arch and the prescribed
equal signs on every torn arch.  Let $\mathcal A(D)$ be the set of
particle-state pairs $(Q_1,Q_2)$ obtained in this way.  The diagram
contributes
\[
 2^{c(D)}\sum_{(Q_1,Q_2)\in\mathcal A(D)}
           e_{Q_1}\otimes e_{Q_2}.
\]
Each permitted endpoint assignment is counted once.  In particular,
a prescribed tear replaces the two ordinary arch assignments by one
new assignment; it is not applied separately to each old coloring.
Summing these contributions over diagrams, marks, and their specified
weights defines the marked vector.

\begin{lemma}[Marked cap vectors]\label{lem:marked-cap}
Start with an ordinary two-replica cap vector.  In each uncolored cap
diagram specify a collection of eligible arches, and retain only diagrams
having at most $M$ eligible arches.  Sum over ordered lists of $k$
distinct eligible arches, with one positive or negative tear prescribed
on each marked arch, and give every list any weight in $[0,1]$ depending
only on the uncolored diagram and its marks.  The resulting vector $V$ satisfies
\begin{equation}\label{eq:marked-cap-norm}
 \|V\|\le M^k\|V_{\rm ordinary}\|.
\end{equation}
Moreover, when such vectors are joined by ordinary double-replica
transfer, their expansion over uncolored configurations has precisely
the ordinary edge multiplicities.  Dividing a contraction by
$Z_{\rm dd}$ therefore gives the expectation of the ratio of modified
color counts to ordinary color counts, with the prescribed mark weights.
The same interpretation holds for an inserted linear operation that
preserves the total seam-edge multiplicity at each site; signs of that
operation are included in the modified count.
\end{lemma}

\begin{proof}
Expand $\|V\|^2$ by pairs of uncolored cap diagrams.  A nonzero gluing
requires their seam multiplicities to agree.  Their two arch pairings
then form alternating cycles on the single seam endpoints.  Every such
cycle has even length.  Before any modification, its opposite-sign
constraints admit exactly two choices.  The tear rules replace some
of these constraints by fixed endpoint signs.  If a cycle has a tear,
those fixed signs force the signs along every remaining segment of
ordinary arches, so there is at most one compatible assignment.
An unmodified cycle still has its two choices.  Thus the number of
assignments cannot increase, although the new assignments need not be
ordinary colorings.  The weights of closed loops internal to either cap
are unchanged.  For each pair of diagrams there are at most $M^k$ lists
on either side of the scalar product.  Termwise comparison with the
ordinary squared norm, followed by summation, proves
\eqref{eq:marked-cap-norm}.

For the second assertion, expand the central transfers as matching
configurations and glue along equal seam states.  Erasing all colors
leaves an ordinary double-dimer configuration: every vertex has degree
two, counting a doubled edge twice, and all single cycles have even
length because the lattice is bipartite.  The modification only changes
the allowed color contractions along those cycles.  An ordinary
uncolored configuration with $r$ single cycles has weight $2^r$, and the
modified contraction supplies its modified color count instead.  This
is exactly the stated ratio after normalization.  A sitewise
multiplicity-preserving insertion changes no part of this uncolored
gluing, so its matrix signs can be incorporated in the same expansion.
\end{proof}

\subsection{Proof of traversal tightness}

\begin{proof}[Proof of Proposition~\ref{prop:rectangle-crossings}]
We use longitudinal coordinates in columns and put
$d=(v-u)/n$.  Choose seams $a,b$ at approximately the one-third and
two-thirds points of $[u,v]$, so each of the three widths is at least
$(v-u)/4$.  In the left cap at $a$, call an arch eligible if it reaches
$x\le u$; in the right cap at $b$, call it eligible if it reaches
$x\ge v$.  Every eligible arch contributes two disjoint traversals of
the corresponding outer subslab.  Hence the probability of more than
$M$ eligible arches on either side is bounded by
\begin{equation}\label{eq:eligible-exception}
 \PP\{N(u,a)\ge2M\}+\PP\{N(b,v)\ge2M\}.
\end{equation}

Suppose now that there are at least $2k$ whole-slab traversals and both
eligible counts are at most $M$.  A loop contributing $2r$ traversals
has cyclic deep-side word $(LR)^r$, where $L$ and $R$ denote visits to
$x\le u$ and $x\ge v$.  Distribute $k$ left/right pairs among the loops,
choosing on each used loop a consecutive block of $s\le r$ pairs.
The selected letters alternate cyclically even across the omitted block:
the last selected $R$ is followed by the first selected $L$.

For each chosen run, take a visit to its deep region and select the
cap arch containing that visit.  This arch is eligible.  The chosen
arches are distinct, because a single cap arch cannot contain a visit
to the opposite deep region and therefore cannot serve two runs
separated by such a visit.  Thus the selected cap arches have precisely
the cyclic left/right order prescribed by the chosen words.

Prescribe positive tears at all $k$ marks in each cap.  Both seam sectors
are now $(p+k,p-k)$.  Equal positive crossing signs describe a source
on a left-cap arch and a sink on a right-cap arch.  The cyclic alternation
of the marks therefore permits a modified matching assignment: continue
the alternating colors between successive marks, switching the endpoint
rule at each mark.  A touched ordinary cycle loses at most its factor
two, while untouched cycles retain it.  There are at most $2k$ touched
cycles, so this marking has modified-to-ordinary color ratio at least
$2^{-2k}$.

Sum over all choices of $k$ eligible marks in each cap.  Every coefficient
is nonnegative.  Lemma~\ref{lem:marked-cap}, followed by
Lemmas~\ref{lem:spectrum} and~\ref{lem:cap-comparison}, bounds the
normalized contraction by
\[
 C M^{2k}\exp(-c d k^2).
\]
Here the constant $C$ is uniform for every cut inside the original slab,
since that slab stays a fixed positive distance from the longitudinal
ends of the rectangle.  Comparing this upper bound to the preceding
pointwise lower bound and using \eqref{eq:eligible-exception} yields
\begin{align}
 \PP\{N(u,v)\ge2k\}
 &\le \PP\{N(u,a)\ge2M\}
       +\PP\{N(b,v)\ge2M\}\notag\\
 &\quad+C\,2^{2k}M^{2k}e^{-c d k^2}.
 \label{eq:crossing-recursion}
\end{align}
Since $N(u,v)$ is even and $n=2p+1$, \eqref{eq:seam-max-crossings}
implies that a nonempty event $N(u,v)\ge2k$ has $k\le p$.
Thus every required sector lies in the range of
Lemma~\ref{lem:spectrum}.

It remains to close the recursion without assuming any prior crossing
bound.  Take an integer $k\ge2$, set $M=k^2$, and apply
\eqref{eq:crossing-recursion} to the two outer slabs.  At depth $j$ the
threshold parameter is
\[
 k_j=k^{2^j},
\]
there are at most $2^j$ slabs, and their widths are at least
$4^{-j}(v-u)$.  The total contribution of the explicit error terms is
at most
\begin{equation}\label{eq:recursion-sum}
 C\sum_{j\ge0}2^j
 \exp\!\left(2k_j\log2+4k_j\log k_j
                      -c d4^{-j}k_j^2\right).
\end{equation}
For all sufficiently large initial $k$, the two positive terms in the
exponent are at most half the negative term, simultaneously for all
$j$: indeed $k_j/(4^j\log k_j)$ has its minimum at $j=0$ once $k$ is
large.  The remaining summands are bounded by
$C2^j\exp(-c' d4^{-j}k^{2^{j+1}})$; their sum tends to zero as
$k\to\infty$.

For a fixed lattice, stop a branch as soon as $2k_j>n$, when its event
is empty by \eqref{eq:seam-max-crossings}.  This happens after
$O(\log\log n)$ levels.  Up to that depth, the smallest available width
is at least $c n/(\log n)^C$ columns for fixed initial $k,d>0$.
Thus all the rounded subdivisions used above are possible once the
mesh is sufficiently small.  No terminal probability remains, and
\eqref{eq:recursion-sum} proves the proposition.  Rotation or reflection
puts either lattice direction in the same convention.
\end{proof}
